\scanpage{118}
Regardons d'abord un sous-fibré vect. de rg 2 quelconque, $H$, et regardons \deleted{\unsure{la}} la forme \deleted{det} quadratique
\[
\det{}_{|H}
\]
On a vu que son discriminant \inserted{$\delta_H$} est aussi égal : celui de la forme bilinéaire \inserted{sym.} $\varphi|H$ $\operatorname{Tr}uv$ (car $H\ni1$). C'est une section de
\[
\det H^{\otimes(\unsure{-|{\deleted{+}}-}2)}\simeq\deleted{\det}(H/k)^{\otimes-2}
\]

Par rapport à une base $1,u$, d'où la base $1\wedge u$ de $\det H\simeq H/k$, et $(1\wedge u)^{\otimes-2}$ de $(\det H)^{\otimes-2}$, on trouve le discriminant
\[
\det\begin{pmatrix}\operatorname{Tr}1&\operatorname{Tr}u\\\operatorname{Tr}u&\operatorname{Tr}u^2\end{pmatrix}=b^2-4c
\]
donc
\[
\delta_H=(b^2-4c)(1\wedge u)^{\otimes-2}
\]

Si on a $H,H'$ tels que $H\cap H'=k$ (sous-vect. de rang 2 en position transverse p.r. à $k.1$), et posant
\[
B=H+H'
\]
comme dessus, on a
\[
B/k\simeq H/k+H'/k
\]
d'où isom (défini au signe près, si on ne précise les conventions)
\[
\det B\simeq\underset{k}{\underline{\det k}}\otimes\det B/k\simeq \deleted{\unsure{\det}}\unsure{H/k\otimes H/k|H'/k\otimes H/k|H/k\otimes H'/k},
\]
Donc \deleted{on peut interpréter}
\[
(\det B)^{\otimes-2}\simeq(H/k)^{\otimes-2}\otimes(H'/k)^{\otimes-2}
\]
% archive sheet number: 117
